\section{Morphisms of finite presentation}
\label{section-finite-presentation}

\noindent
Recall that a ring map $R \to A$ is of finite presentation if
$A$ is isomorphic to $R[x_1, \ldots, x_n]/(f_1, \ldots, f_m)$ as
an $R$-algebra for some $n, m$ and some polynomials $f_j$, see
Algebra, Definition \ref{algebra-definition-finite-type}.

\begin{definition}
\label{definition-finite-presentation}
Let $f : X \to S$ be a morphism of schemes.
\begin{enumerate}
\item We say that $f$ is of {\it finite presentation at $x \in X$} if
there exists an affine open neighbourhood $\Spec(A) = U \subset X$
of $x$ and affine open $\Spec(R) = V \subset S$
with $f(U) \subset V$ such that the induced ring map
$R \to A$ is of finite presentation.
\item We say that $f$ is {\it locally of finite presentation} if it is
of finite presentation at every point of $X$.
\item We say that $f$ is of {\it finite presentation} if it is locally of
finite presentation, quasi-compact and quasi-separated.
\end{enumerate}
\end{definition}

\noindent
Note that a morphism of finite presentation is {\bf not} just a quasi-compact
morphism which is locally of finite presentation.
Later we will characterize morphisms which are
locally of finite presentation as those morphisms such that
$$
\colim \Mor_S(T_i, X) = \Mor_S(\lim T_i, X)
$$
for any directed system of affine schemes $T_i$ over $S$. See
Limits,
Proposition \ref{limits-proposition-characterize-locally-finite-presentation}.
In Limits, Section \ref{limits-section-descending-relative} we show
that, if $S = \lim_i S_i$ is a limit of affine schemes,
any scheme $X$ of finite presentation over $S$ descends to a scheme
$X_i$ over $S_i$ for some $i$.

\begin{lemma}
\label{lemma-locally-finite-presentation-characterize}
Let $f : X \to S$ be a morphism of schemes.
The following are equivalent
\begin{enumerate}
\item The morphism $f$ is locally of finite presentation.
\item For every affine opens $U \subset X$, $V \subset S$
with $f(U) \subset V$ the ring map
$\mathcal{O}_S(V) \to \mathcal{O}_X(U)$ is of finite presentation.
\item There exist an open covering $S = \bigcup_{j \in J} V_j$
and open coverings $f^{-1}(V_j) = \bigcup_{i \in I_j} U_i$ such
that each of the morphisms $U_i \to V_j$, $j\in J, i\in I_j$
is locally of finite presentation.
\item There exist an affine open covering $S = \bigcup_{j \in J} V_j$
and affine open coverings $f^{-1}(V_j) = \bigcup_{i \in I_j} U_i$ such
that the ring map $\mathcal{O}_S(V_j) \to \mathcal{O}_X(U_i)$ is
of finite presentation, for all $j\in J, i\in I_j$.
\end{enumerate}
Moreover, if $f$ is locally of finite presentation then for
any open subschemes $U \subset X$, $V \subset S$ with $f(U) \subset V$
the restriction $f|_U : U \to V$ is locally of finite presentation.
\end{lemma}

\begin{proof}
This follows from Lemma \ref{lemma-locally-P-characterize} if we show that
the property ``$R \to A$ is of finite presentation'' is local.
We check conditions (a), (b) and (c) of Definition
\ref{definition-property-local}.
By Algebra, Lemma \ref{algebra-lemma-base-change-finiteness}
being of finite presentation is stable under base change and hence
we conclude (a) holds. By Algebra, Lemma
\ref{algebra-lemma-compose-finite-type} being of finite presentation
is stable under composition and trivially for any ring
$R$ the ring map $R \to R_f$ is of finite presentation.
We conclude (b) holds. Finally, property (c) is true
according to Algebra, Lemma \ref{algebra-lemma-cover-upstairs}.
\end{proof}

\begin{lemma}
\label{lemma-composition-finite-presentation}
The composition of two morphisms which are locally of finite presentation is
locally of finite presentation.
The same is true for morphisms of finite presentation.
\end{lemma}

\begin{proof}
In the proof of Lemma \ref{lemma-locally-finite-presentation-characterize}
we saw that being of finite presentation is a local property of ring maps.
Hence the first statement of the lemma follows from
Lemma \ref{lemma-composition-type-P} combined
with the fact that being of finite presentation
is a property of ring maps that is
stable under composition, see
Algebra, Lemma \ref{algebra-lemma-compose-finite-type}.
By the above and the fact that compositions of
quasi-compact, quasi-separated morphisms are
quasi-compact and quasi-separated, see
Schemes, Lemmas \ref{schemes-lemma-composition-quasi-compact}
and \ref{schemes-lemma-separated-permanence}
we see that the composition of morphisms of finite presentation is
of finite presentation.
\end{proof}

\begin{lemma}
\label{lemma-base-change-finite-presentation}
The base change of a morphism which is locally of finite presentation
is locally of finite presentation. The same is true for morphisms of
finite presentation.
\end{lemma}

\begin{proof}
In the proof of Lemma \ref{lemma-locally-finite-presentation-characterize}
we saw that being of finite presentation is a local property of ring maps.
Hence the first statement of the lemma follows from
Lemma \ref{lemma-composition-type-P} combined
with the fact that being of finite presentation
is a property of ring maps that is
stable under base change, see
Algebra, Lemma \ref{algebra-lemma-base-change-finiteness}.
By the above and the fact that a base change of a
quasi-compact, quasi-separated morphism is quasi-compact
and quasi-separated, see
Schemes, Lemmas \ref{schemes-lemma-quasi-compact-preserved-base-change}
and \ref{schemes-lemma-separated-permanence}
we see that the base change of a morphism of finite presentation is
a morphism of finite presentation.
\end{proof}

\begin{lemma}
\label{lemma-open-immersion-locally-finite-presentation}
Any open immersion is locally of finite presentation.
\end{lemma}

\begin{proof}
This is true because an open immersion is a local isomorphism.
\end{proof}

\begin{lemma}
\label{lemma-quasi-compact-open-immersion-finite-presentation}
Any open immersion is of finite presentation if and only if
it is quasi-compact.
\end{lemma}

\begin{proof}
We have seen (Lemma \ref{lemma-open-immersion-locally-finite-presentation})
that an open immersion is locally of finite presentation.
We have seen (Schemes, Lemma \ref{schemes-lemma-immersions-monomorphisms})
that an immersion is separated and hence quasi-separated. From this
and Definition \ref{definition-finite-presentation} the lemma follows.
\end{proof}

\begin{lemma}
\label{lemma-closed-immersion-finite-presentation}
\begin{slogan}
Closed immersions of finite presentation correspond
to quasi-coherent sheaves of ideals of finite type.
\end{slogan}
A closed immersion $i : Z \to X$ is of finite presentation if and only if
the associated quasi-coherent sheaf of ideals
$\mathcal{I} = \Ker(\mathcal{O}_X \to i_*\mathcal{O}_Z)$
is of finite type (as an $\mathcal{O}_X$-module).
\end{lemma}

\begin{proof}
On any affine open $\Spec(R) \subset X$ we have
$i^{-1}(\Spec(R)) = \Spec(R/I)$ and
$\mathcal{I} = \widetilde{I}$. Moreover, $\mathcal{I}$
is of finite type if and only if $I$ is a finite $R$-module
for every such affine open (see
Properties, Lemma \ref{properties-lemma-finite-type-module}).
And $R/I$ is of finite presentation
over $R$ if and only if $I$ is a finite $R$-module. Hence we win.
\end{proof}

\begin{lemma}
\label{lemma-finite-presentation-finite-type}
A morphism which is locally of finite presentation is locally of finite type.
A morphism of finite presentation is of finite type.
\end{lemma}

\begin{proof}
Omitted.
\end{proof}

\begin{lemma}
\label{lemma-noetherian-finite-type-finite-presentation}
\begin{slogan}
Over a locally Noetherian base, finite type is finite presentation.
\end{slogan}
Let $f : X \to S$ be a morphism.
\begin{enumerate}
\item If $S$ is locally Noetherian and $f$ locally of finite type
then $f$ is locally of finite presentation.
\item If $S$ is locally Noetherian and $f$ of finite type
then $f$ is of finite presentation.
\end{enumerate}
\end{lemma}

\begin{proof}
The first statement follows from the fact that a ring
of finite type over a Noetherian ring is of finite presentation, see Algebra,
Lemma \ref{algebra-lemma-Noetherian-finite-type-is-finite-presentation}.
Suppose that $f$ is of finite type and $S$ is locally Noetherian.
Then $f$ is quasi-compact and locally of finite presentation by (1).
Hence it suffices to prove that $f$ is quasi-separated.
This follows from Lemma \ref{lemma-finite-type-Noetherian-quasi-separated}
(and Lemma \ref{lemma-finite-presentation-finite-type}).
\end{proof}

\begin{lemma}
\label{lemma-finite-presentation-quasi-compact-quasi-separated}
Let $S$ be a scheme which is quasi-compact and quasi-separated.
If $X$ is of finite presentation over $S$, then $X$ is quasi-compact
and quasi-separated.
\end{lemma}

\begin{proof}
Omitted.
\end{proof}

\begin{lemma}
\label{lemma-finite-presentation-permanence}
Let $f : X \to Y$ be a morphism of schemes over $S$.
\begin{enumerate}
\item If $X$ is locally of finite presentation over $S$ and
$Y$ is locally of finite type over $S$, then $f$ is locally
of finite presentation.
\item If $X$ is of finite presentation over $S$ and $Y$ is quasi-separated
and locally of finite type over $S$, then $f$ is of finite presentation.
\end{enumerate}
\end{lemma}

\begin{proof}
Proof of (1). Via Lemma \ref{lemma-locally-finite-presentation-characterize}
this translates into the following algebra
fact: Given ring maps $A \to B \to C$ such that $A \to C$ is
of finite presentation and $A \to B$ is of finite type,
then $B \to C$ is of finite presentation. See
Algebra, Lemma \ref{algebra-lemma-compose-finite-type}.

\medskip\noindent
Part (2) follows from (1) and
Schemes, Lemmas \ref{schemes-lemma-compose-after-separated} and
\ref{schemes-lemma-quasi-compact-permanence}.
\end{proof}

\begin{lemma}
\label{lemma-diagonal-morphism-finite-type}
Let $f : X \to Y$ be a morphism of schemes with diagonal
$\Delta : X \to X \times_Y X$. If $f$ is locally of finite type
then $\Delta$ is locally of finite presentation. If $f$ is quasi-separated
and locally of finite type, then $\Delta$ is of finite presentation.
\end{lemma}

\begin{proof}
Note that $\Delta$ is a morphism of schemes over $X$ (via the second
projection $X \times_Y X \to X$). Assume $f$ is locally of finite type.
Note that $X$ is of finite presentation over $X$ and $X \times_Y X$ is
locally of finite type over $X$ (by Lemma \ref{lemma-base-change-finite-type}).
Thus the first statement holds by
Lemma \ref{lemma-finite-presentation-permanence}.
The second statement follows from the first, the definitions, and
the fact that a diagonal morphism is a monomorphism, hence separated
(Schemes, Lemma \ref{schemes-lemma-monomorphism-separated}).
\end{proof}








